\documentclass[12pt,letterpaper]{article}

% Librerías a utilizar
\usepackage[utf8]{inputenc}	% Codificación 
\usepackage[spanish]{babel}	% Idioma
\usepackage{natbib}			% Bibliografía
\usepackage{graphicx}		% Imagenes
\usepackage{indentfirst}		% Sangría
\usepackage{amsmath, amsfonts, amssymb}	% Figuras matemáticas
\usepackage{url}							% URL
\usepackage{subfigure}
\usepackage{svg}
\usepackage{algpseudocode,algorithm}
\floatname{algorithm}{Algoritmo}
\usepackage[toc,titletoc,title]{appendix}
\renewcommand{\appendixtocname}{Anexos}
\renewcommand{\appendixpagename}{Anexos}
\usepackage{tabularx}
\usepackage{pdflscape}
\usepackage{xltabular}
\usepackage{multirow}
\usepackage{booktabs}
\usepackage[table]{xcolor}
\usepackage{colortbl}
\usepackage{array}\setlength{\extrarowheight}{1pt}
\usepackage{rotating}
\usepackage{stfloats}
\usepackage{placeins}
\usepackage{float}
\usepackage{caption}
\usepackage{enumitem}

% Enlaces cliqueables en el PDF (citas, índice, figuras, etc.)
\usepackage[hidelinks]{hyperref}
\newcommand{\appendixnumberline}[1]{\makebox[5.5em][l]{#1}}

\tolerance=1
\emergencystretch=\maxdimen
\hyphenpenalty=10000
\hbadness=10000

\renewcommand{\theenumi}{\alph{enumi}} %Enumerar Letras minúsculas 

\usepackage[left=3cm,right=2cm,top=2cm,bottom=3cm]{geometry}

\setlength{\parindent}{2cm}	% Sangría en los párrafos
\renewcommand{\baselinestretch}{1.5}	% Interlineado

% Renombrar ciertos títulos del texto
\renewcommand\spanishcontentsname{Tabla de contenidos de Anexos}
\renewcommand\spanishlisttablename{Índice de tablas}
\renewcommand\spanishrefname{Referencias}
\renewcommand\spanishtablename{Tabla}

\begin{document}

\newpage
\thispagestyle{empty}

\begin{center}
	{\bf \mbox{\large UNIVERSIDAD DE SANTIAGO DE CHILE}}\\
	{\bf \mbox{FACULTAD DE INGENIERÍA}}\\
	{\bf \mbox{DEPARTAMENTO DE INGENIERÍA INFORMÁTICA}}\\
\end{center}

\vspace{4cm}
\begin{center}
	\Large
	\textbf{Anexos Digitales del Trabajo de Título}
	
	\vspace{1cm}
	\large
	\textbf{Metodología basada en procesamiento de lenguaje natural para analizar la alineación curricular de Ingeniería Informática en Chile con estándares ACM/IEEE y necesidades de la industria}
\end{center}

\vspace*{2.5cm}
\begin{flushright}
	\begin{tabular}[t]{l l}
		Estudiante:
		               & Fabián Ibarra Norambuena \\
		Profesor guía: & Manuel Villalobos Cid    \\
	\end{tabular}
\end{flushright}
\begin{center}
	\vspace{2.5cm}
	Santiago - Chile\\
	2026
\end{center}

\newpage
\pagenumbering{roman}
\tableofcontents

\newpage
\pagenumbering{arabic}
\setcounter{page}{1}

\section*{Anexos}
\addcontentsline{toc}{section}{Anexos}
\setcounter{subsection}{0}
\renewcommand{\thesubsection}{Anexo \Alph{subsection}}
\renewcommand{\theHsubsection}{Anexo.\arabic{subsection}}
\addtocontents{toc}{\string\renewcommand{\string\numberline}{\string\appendixnumberline}}

\subsection{Glosario}

\begin{itemize}
    \item \textbf{Alineación Curricular}: Grado de correspondencia, congruencia y consistencia entre los componentes de un programa educativo (tales como perfiles de egreso, planes de estudio y descriptores de asignaturas) con respecto a directrices internacionales y los requerimientos del entorno laboral.
    
    \item \textbf{Clustering (Agrupamiento)}: Técnica de aprendizaje no supervisado que consiste en agrupar un conjunto de objetos de tal manera que los miembros del mismo grupo (\textit{clúster}) sean más similares entre sí que con los de otros grupos, basándose en alguna métrica de distancia o similitud geométrica.
    
    \item \textbf{Coeficiente de Silueta (\textit{Silhouette Coefficient})}: Métrica de validación interna de algoritmos de \textit{clustering} que evalúa qué tan bien agrupado está un elemento. Mide la cercanía de un objeto a su propio clúster (cohesión) frente a su cercanía a otros clústeres (separación), variando entre -1 y 1.
    
    \item \textbf{Corpus}: Conjunto estructurado y homogéneo de textos escritos, recopilado de manera sistemática para ser utilizado como base de análisis lingüístico o procesamiento computacional.
    
    \item \textbf{DBSCAN (\textit{Density-Based Spatial Clustering of Applications with Noise})}: Algoritmo de agrupamiento no supervisado basado en densidad que agrupa puntos cercanos y etiqueta como valores atípicos (\textit{outliers}) aquellos que se encuentran aislados en regiones de baja densidad.
    
    \item \textbf{Embeddings (Incrustaciones Vectoriales)}: Representaciones matemáticas de textos (palabras, frases o documentos completos) en forma de vectores densos en un espacio continuo de baja dimensionalidad, donde la cercanía geométrica de los vectores representa similitud semántica.
    
    \item \textbf{Índice de Davies-Bouldin}: Métrica de evaluación intrínseca de algoritmos de agrupamiento que calcula la similitud promedio entre cada clúster y su clúster más similar. Valores más bajos indican agrupamientos con clústeres más compactos y mejor separados.
    
    \item \textbf{K-Means}: Algoritmo clásico de agrupamiento no supervisado que particiona un conjunto de datos en $K$ clústeres distintos y disjuntos, asignando cada elemento al clúster cuyo centroide se encuentra a menor distancia geométrica (generalmente distancia euclidiana).
    
    \item \textbf{Lematización}: Proceso de preprocesamiento de lenguaje natural que consiste en reducir las palabras de un texto a su forma básica o raíz morfológica canonizada (lema), como infinitivos para verbos y masculino singular para adjetivos y sustantivos.
    
    \item \textbf{Procesamiento de Lenguaje Natural (PLN / NLP)}: Subcampo de la inteligencia artificial y la lingüística computacional enfocado en el diseño e implementación de algoritmos que permiten a las computadoras procesar, entender y analizar textos en lenguaje humano.
    
    \item \textbf{Sentence-Transformers (SBERT)}: Framework de aprendizaje profundo basado en redes neuronales siamesas, utilizado para adaptar modelos de lenguaje tipo BERT para generar embeddings de oraciones y textos completos de alta calidad semántica.
    
    \item \textbf{Similitud del Coseno}: Métrica matemática utilizada para cuantificar la similitud entre dos vectores en un espacio multidimensional, midiendo el coseno del ángulo formado entre ellos, independientemente de su magnitud.
    
    \item \textbf{Stop Words (Palabras Vacías)}: Términos que carecen de significado informativo relevante por sí mismos (como artículos, preposiciones y conjunciones, o en el contexto académico, términos como ``asignatura'', ``horas'' o ``carrera''), que son filtrados durante el preprocesamiento de texto.
    
    \item \textbf{TF-IDF (\textit{Term Frequency - Inverse Document Frequency})}: Técnica estadística clásica de representación de textos que evalúa la importancia de una palabra en un documento en relación con un corpus completo, ponderándola según su frecuencia local y su rareza global.
    
    \item \textbf{t-SNE (\textit{t-Distributed Stochastic Neighbor Embedding})}: Algoritmo no lineal de reducción de dimensionalidad utilizado principalmente para la visualización de datos de alta dimensionalidad en espacios de dos o tres dimensiones.
    
    \item \textbf{UMAP (\textit{Uniform Manifold Approximation and Projection})}: Algoritmo moderno de reducción de dimensionalidad no lineal fundamentado en la geometría de variedades, eficiente computacionalmente y robusto para preservar la estructura local y global de los datos.
\end{itemize}

%Results from the state-of-the-art strategies
%\setcounter{page}{1}
\subsection{Resumen comparativo de trabajos relacionados sobre alineación curricular automatizada.}
\begin{landscape}
\footnotesize
\renewcommand{\baselinestretch}{1.0}\selectfont
\renewcommand{\arraystretch}{1.2}
\begin{xltabular}{\linewidth}{p{2.2cm} X X X X}
    \caption{Análisis comparativo de objetivos, metodologías, resultados y limitaciones de la literatura de mapeo curricular.} \label{tab:estado_arte} \\
    \toprule
    \textbf{Referencia} &
    \textbf{Objetivo}   &
    \textbf{Metodología} &
    \textbf{Resultados principales} &
    \textbf{Limitaciones identificadas} \\
    \midrule
    \endhead

    \citep{russell_e_mapping_2024} & 
    Automatizar el mapeo de planes de estudio hacia habilidades y ocupaciones de la industria utilizando descripciones de cursos. & 
    Procesamiento de Lenguaje Natural (NLP) usando SkillNER, extracción de características con ponderación TF-IDF y similitud del coseno para emparejar con O*NET. & 
    Extracción promedio de 5.3 a 10.4 habilidades por curso. Fuerte precisión (precision@5 del 60\% al 100\%). Fuerte señal de habilidades incluso en textos breves. & 
    Evaluación limitada a un campo (ciencia de datos) y dos instituciones. Falsos positivos por coincidencia léxica. Falta información sobre la profundidad de la habilidad enseñada. \\
    \addlinespace
    
    \citep{cristea_course_2025} & 
    Evaluar el rendimiento de LLM frente a métodos NLP tradicionales en la extracción de habilidades para la analítica curricular. & 
    Comparación de estrategias: Generación Aumentada por Recuperación (RAG), prompting Zero-shot (GPT, Llama, Gemini, Claude), modelos basados en embeddings (SBERT) y TF-IDF en 400 documentos. & 
    RAG fue la estrategia con mejor desempeño en todos los documentos. El prompting Zero-shot tuvo peor rendimiento que métodos NLP tradicionales en la mayoría de los casos. & 
    El rendimiento varía significativamente según el modelo seleccionado y las estrategias de prompting. Requiere diseño y ajuste cuidado antes de implementaciones a gran escala. \\
    \addlinespace
    
    \citep{cristea_artificial_2025} & 
    Evaluar el rol de la IA Generativa (LLMs) como co-revisor curricular para generar mapeos ponderados de habilidades en evaluaciones. & 
    Uso de GPT-4.0 frente a métodos ML (BERT, TF-IDF) para mapear 7 habilidades empresariales. Uso del alfa de Krippendorff para comparar resultados con el juicio de expertos. & 
    Los mapeos generados por LLM se alinearon más con el juicio experto (alfa 0.76) que los métodos ML (0.65). Alta precisión en competencias como Comunicación Oral y Trabajo en Equipo. & 
    Limitado a un solo LLM y programa de estudios. Desempeño débil en habilidades que se infieren del contexto (ej. Comunicación escrita) y no se detallan en el texto. \\
    \addlinespace
    
    \citep{fan_identifying_2023} & 
    Identificar capacidades digitales evaluadas en rúbricas de cursos universitarios de manera automatizada utilizando Machine Learning. & 
    Extracción de características mediante TF-IDF. Comparación de clasificadores: Redes Neuronales (ANN), Máquinas de Vectores de Soporte (SVM), Random Forest y XGBoost. & 
    SVM obtuvo el mejor rendimiento (exactitud 0.8535 y F1-score 0.8338). El rendimiento mejora y luego decae al aumentar las características extraídas. & 
    Dependencia de la intervención humana para etiquetar los datos de entrenamiento (solo 1.735 muestras etiquetadas). Los algoritmos probados no pueden utilizar pesos difusos para mitigar ruido. \\
    \addlinespace
    
    \citep{duarte_broader_2023} & 
    Generar representaciones visuales comprensibles de los planes de estudio para reducir brechas de comunicación entre los actores. & 
    Clasificación usando índices de Wikipedia, NLP (tokenización y normalización), y visualización mediante ideogramas circulares y grafos de red. & 
    El mapeo apoyó con evidencia visual las debilidades del programa. Las representaciones cuantitativas empoderan el entendimiento entre diferentes grupos de interés. & 
    Las representaciones agrupan conceptos pero no trazan rutas de aprendizaje precisas o de precedencia detalladas, requiriendo revisión y otras herramientas para detalles exactos. \\
    \addlinespace
    
    \citep{miura_automated_2021} & 
    Mapear automáticamente programas educativos de ciencias de la computación frente a estándares de aprendizaje (ITSS+) para reducir costos manuales. & 
    Representación de texto usando Bag-of-Words (BoW) y Word Embeddings (Word2Vec). Emparejamiento usando similitud del coseno y algoritmos de aprendizaje supervisado. & 
    El aprendizaje supervisado con Word Embeddings alcanzó la mayor medida F (0.571), superando ampliamente a enfoques tradicionales como BoW (0.333) al capturar semántica abstracta. & 
    Requiere aumentar los datos de entrenamiento y cuantificar más correctamente características ponderando partes importantes del material didáctico. \\
    \addlinespace
    
    \citep{constantinov_preliminary_2015} & 
    Ajustar un currículo de ingeniería extrayendo las demandas del mercado laboral a través de redes profesionales. & 
    Pipeline de datos para Análisis de Redes Profesionales (PNA) con perfiles de LinkedIn. Extracción de habilidades en sílabos mediante TF-IDF y diccionarios personalizados. & 
    El sistema identificó habilidades del plan de estudios obsoletas en la industria y habilidades muy demandadas en el trabajo que el programa académico no estaba cubriendo. & 
    Metodología basada en diccionarios personalizados estáticos y extracción mediante palabras clave sin abordajes semánticos complejos o contextuales de modelado de lenguaje. \\
  
    \bottomrule
\end{xltabular}
\end{landscape}




% Anexo 3
\subsection{Descripción del corpus de datos}
\label{app:corpus}
 \begin{table}[H]
        \centering
        \small
        % interlineado simple dentro de la tabla (1.0)
        \renewcommand{\baselinestretch}{1.0}\selectfont
        \caption{Planes de estudio y asignaturas recopiladas por universidad.}
        \label{tab:corpus_universidades}
        \captionsetup{skip=8pt}                                                                                              
                                                                                                                             
        % Definimos el contador de filas (se resetea al inicio de la tabla)
        \newcounter{contadorFilas}
        \setcounter{contadorFilas}{0}
        
        % Comando para terminar la fila incrementando el contador y dibujando una línea horizontal
        \newcommand{\finFila}{%
          \stepcounter{contadorFilas}%
          \\ \hline%
        }
        
        % Definición de columnas con alineación a la izquierda (raggedright) para evitar espacios gigantes
        \begin{tabularx}{\textwidth}{>{\raggedright\arraybackslash}X >{\raggedright\arraybackslash}X c}                                                                              
            \toprule                                                                                                         
            \textbf{Universidad} &                                                                                           
            \textbf{Carrera / Programa} &                                                                                    
            \textbf{N\textsuperscript{o} Asignaturas} \\                                                                                    
            \midrule                                                                                                         
                                                                                                                             
            % Fila 1                                                                                              
            Universidad de Santiago de Chile &                                                                               
            Ingeniería de Ejecución en Computación e Informática &                                                                                
            43 \finFila                                                                                       
                                                                                                                             
            % Fila 2                                                                                              
            Universidad de Santiago de Chile &                                                                                           
            Ingeniería Civil en Informática &                                                                                
            58 \finFila
            
            % Fila 3 
            Universidad de Chile & 
            Ingeniería Civil en Computación & 
            54 \finFila
            
            % Fila 4
            Universidad Andrés bello & 
            Ingeniería en Computación e informática & 
            41 \finFila
            
            % Fila 5
            Universidad de Magallanes & 
            Ingeniería en Computación e Informática & 
            41 \finFila

            % Fila 6
            Universidad Tecnológica Metropolitana & 
            Ingeniería en Informática & 
            44 \finFila

            % Fila 7
            Universidad de Concepción & 
            Ingeniería Civil Informática & 
            54 \finFila

            % Fila 8
            Universidad Técnica Federico Santa María & 
            Ingeniería en Informática & 
            42 \finFila

            % Fila 9
            Pontificia Universidad Catolica de Valparaiso  & 
            Ingeniería Informática & 
            44 \finFila

            % Fila 10
            DuocUC  & 
            Ingeniería Informática & 
            33 \finFila

            % Fila 11
            Universidad San Sebastián  & 
            Ingeniería Civil en Informática & 
            50 \finFila

            % Fila 12
            Universidad Santo Tomás   & 
            Ingeniería Civil en Informatica y Sistemas Inteligentes  & 
            60 \finFila

            % Fila 13
            Universidad de la Frontera   & 
            Ingeniería Civil Informática  & 
            57 \finFila

            % Fila 14
            Universidad del Bío-Bío   & 
            Ingeniería de Ejecución en Computación e Informática  & 
            46 \finFila

            % Fila 15
            Universidad de Talca   & 
            Ingeniería Civil en Computación  & 
            43 \finFila

            % Fila 16
            Universidad Autonoma  & 
            Ingeniería en Informática  & 
            32 \finFila
            
            \multicolumn{2}{l}{\textbf{Total del Corpus (16 Programas / Planes de Estudio)}} & \textbf{747} \\
            \bottomrule
        \end{tabularx}
    \end{table}

% Anexo 4
\subsection{Protocolo y caracterización de la muestra de datos de LinkedIn}
\label{app:linkedin_sample}

Con el fin de garantizar la representatividad, validez y reproducibilidad del análisis de pertinencia laboral, se define un protocolo formal para la selección y curaduría de la muestra de ofertas de empleo provenientes de la plataforma LinkedIn:

\begin{enumerate}[label=\alph*.]
    \item \textbf{Tamaño y estratificación muestral objetivo ($N$)}: Se establece una muestra total de $N = 200$ anuncios de empleo curados y estructurados recopilados en Chile (2024–2026). Para capturar adecuadamente la diversidad de la demanda laboral local, la muestra se divide formalmente en dos estratos:
    \begin{itemize}
        \item \textbf{Sector Tecnológico (Mundo TI):} Empresas cuyo giro principal es el desarrollo de software, servicios de consultoría TI, telecomunicaciones y startups tecnológicas.
        \item \textbf{Sector No-TI (Industria Usuaria):} Empresas de sectores tradicionales (banca e instituciones financieras, minería, retail, salud, energía y servicios públicos) que requieren perfiles informáticos para sus departamentos de tecnología interna y transformación digital.
    \end{itemize}
    \item \textbf{Ventana temporal de captura}: Avisos de trabajo publicados activamente entre enero de 2024 y junio de 2026.
    \item \textbf{Categorías de cargos consideradas}: La selección abarca los perfiles profesionales más representativos del mercado de TI en Chile:
    \begin{itemize}
        \item Desarrollador de Software / Ingeniero de Software (\textit{Backend}, \textit{Frontend}, \textit{Fullstack}).
        \item Ingeniero / Arquitecto de Datos (\textit{Data Engineer / Architect}).
        \item Científico de Datos / Analista de Inteligencia Artificial (\textit{Data Scientist / AI Specialist}).
        \item Ingeniero de Infraestructura y Operaciones (\textit{DevOps / Cloud Engineer}).
        \item Ingeniero en Ciberseguridad y Redes.
        \item Arquitecto de Software y Sistemas.
    \end{itemize}
    \item \textbf{Criterios de inclusión y exclusión}:
    \begin{itemize}
        \item \textit{Inclusión}: Ofertas laborales localizadas en Chile (presencial, híbrido o remoto), con descripción detallada de funciones, requisitos técnicos explícitos y nivel de experiencia (junior, semi-senior, senior).
        \item \textit{Exclusión}: Avisos duplicados, ofertas de pasantías o prácticas profesionales no remuneradas, solicitudes con descripciones genéricas sin requerimientos tecnológicos concretos, y vacantes orientadas puramente a ventas o gestión comercial sin exigencia técnica.
    \end{itemize}
    \item \textbf{Protocolo de curaduría y lematización}: Extracción de textos mediante estructuración de secciones de ``Requisitos'' y ``Responsabilidades'', seguida por filtrado de palabras vacías laborales (ej. ``empresa'', ``beneficios'', ``postular'') y extracción de entidades tecnológicas (lenguajes, frameworks, metodologías).
    \item \textbf{Cumplimiento de Términos de Servicio (ToS)}: La captura de datos se realiza en estricto cumplimiento con los Términos de Servicio de LinkedIn, utilizando información de acceso público de forma manual o mediante APIs de consulta de vacantes públicas para fines exclusivamente de investigación académica sin fines de lucro, sin scraping masivo ni almacenamiento de datos personales identificables de reclutadores o postulantes.
\end{enumerate}

% Anexo 5
\subsection{Flujo de procesamiento metodológico}
\label{app:flujo_metodologico}

\begin{table}[H]
    \centering
    \small
    \renewcommand{\baselinestretch}{1.0}\selectfont
    \caption{Detalle de entradas, procesos y salidas del flujo metodológico mínimo.}
    \label{tab:flujo_metodologico_detalle}
    \captionsetup{skip=8pt}
    \begin{tabularx}{\textwidth}{l >{\raggedright\arraybackslash}X >{\raggedright\arraybackslash}X}
        \toprule
        \textbf{Fase} & \textbf{Entrada / Fuentes} & \textbf{Resultado / Salida} \\
        \midrule
        1. Recolección & Mallas, PDF de asignaturas, estándares, LinkedIn & Documentos crudos en base de datos \\
        2. Limpieza & Textos crudos extraídos & Textos planos limpios en minúsculas \\
        3. Normalización & Nombres de asignaturas heterogéneos & Nombres estandarizados y homologados \\
        4. Construcción Corpus & Textos normalizados y filtrados & Base de datos del corpus estructurado \\
        5. Rep. Vectorial & Corpus estructurado & Vectores densos (embeddings multilingües) \\
        6. Reducción Dim. & Vectores de alta dimensionalidad & Vectores en espacio de menor dimensión \\
        7. Clustering & Vectores de asignaturas/perfiles & Clústeres (K-Means, Jerárquico) y Outliers (DBSCAN) \\
        8. Similitud & Vectores y centroides de grupos & Índices de brecha y alineación semántica \\
        9. Visualización & Métricas de alineación y clústeres & Gráficos interactivos y reportes visuales \\
        \bottomrule
    \end{tabularx}
\end{table}

% Anexo 6
\subsection{Planificación temporal del proyecto (Carta Gantt)}
\label{app:gantt}

\begin{table}[H]
    \centering
    \small
    \renewcommand{\baselinestretch}{1.0}\selectfont
    \caption{Carta Gantt con la distribución temporal de las actividades (en semanas).}
    \label{tab:carta_gantt}
    \captionsetup{skip=8pt}
    \begin{tabularx}{\textwidth}{>{\raggedright\arraybackslash}p{6cm} c c c c c c c c}
        \toprule
        \textbf{Actividad / Semana} & \textbf{1-2} & \textbf{3-4} & \textbf{5-6} & \textbf{7-8} & \textbf{9-10} & \textbf{11-12} & \textbf{13-14} & \textbf{15-16} \\
        \midrule
        Fase 1: Recolección y comprensión & \cellcolor{gray!40}X & \cellcolor{gray!40}X & & & & & & \\
        Fase 2: Preparación de datos & & \cellcolor{gray!40}X & \cellcolor{gray!40}X & & & & & \\
        Fase 3: Modelamiento y embeddings & & & \cellcolor{gray!40}X & \cellcolor{gray!40}X & \cellcolor{gray!40}X & & & \\
        Fase 4: Evaluación y métricas & & & & & \cellcolor{gray!40}X & \cellcolor{gray!40}X & & \\
        Fase 5: Despliegue (Dashboard) & & & & & & \cellcolor{gray!40}X & \cellcolor{gray!40}X & \\
        Redacción informe y validación & \cellcolor{gray!20}X & \cellcolor{gray!20}X & \cellcolor{gray!20}X & \cellcolor{gray!20}X & \cellcolor{gray!20}X & \cellcolor{gray!20}X & \cellcolor{gray!40}X & \cellcolor{gray!40}X \\
        \bottomrule
    \end{tabularx}
\end{table}



\newpage
\addcontentsline{toc}{section}{Referencias de Anexos}
\bibliographystyle{apalike}
\bibliography{References/references.bib}

\end{document}
